%% shortname: Uncompensated Law of Demand and the Representative Consumer
\documentclass[11pt]{article}
\usepackage{microstyle}

\begin{document}

\lessonheader{13}{Chapter 4: aggregation across consumers --- the uncompensated law of demand, aggregate WARP and representative consumers}{}


\section{Question 2: when does aggregate demand satisfy WARP?}

\runin{Graphical example.} By example (Lesson 11): not always, even when every
individual demand satisfies WARP. But we know that WARP is the same as the
compensated law of demand (CLD).

$\Rightarrow$ Let us look for an alternative to the CLD.

\topic{The uncompensated law of demand (ULD)}
A law of demand for \emph{uncompensated} price changes: we hold income $w$ fixed.
The demand function $x_i(p,w_i)$ satisfies the ULD if
\[ (p'-p)\cdot\bigl(x_i(p',w_i)-x_i(p,w_i)\bigr)\le 0
   \qquad\text{for any } p,\ p' \text{ and } w_i, \]
with strict inequality whenever $x_i(p',w_i)\ne x_i(p,w_i)$.

The ULD rules out some preferences: it rules out Giffen goods (where consumption
$\uparrow$ when price $\uparrow$).

However, homothetic preferences satisfy the ULD.

\topic{The ULD aggregates}
Consider the wealth distribution rule
\[ w_i(p,w)=\alpha_i w,\qquad \alpha_i>0,\quad \sum_{i=1}^{I}\alpha_i=1. \]
Then the ULD aggregates: if it holds for all agents, then it holds in the
aggregate. (Simply sum the individual ULDs: with $w_i=\alpha_i w$ fixed,
$\sum_i (p'-p)\cdot\bigl(x_i(p',\alpha_i w)-x_i(p,\alpha_i w)\bigr)\le 0$.)


\section{The aggregate ULD implies aggregate WARP}

In fact, the aggregate ULD implies that aggregate demand satisfies WARP.

\runin{Proof.} The proof relies in part on individual homogeneity of degree 0 ---
it is useful in proofs.

Take $(p,w)$ and $(p',w')$ with $x(p,w)\ne x(p',w')$ and
\[ p\cdot x(p',w')\le w. \]
That is, $x(p,w)$ is revealed preferred to $x(p',w')$. We must show
$p'\cdot x(p,w)>w'$.

Aggregate demand is homogeneous of degree zero because individual demands are:
\[ x(p',w')=x\Bigl(\frac{w}{w'}p',\ \frac{w}{w'}w'\Bigr)=x\Bigl(\frac{w}{w'}p',\ w\Bigr), \]
implying, by the revealed-preference assumption and Walras' law,
\[ p\cdot x\Bigl(\frac{w}{w'}p',w\Bigr)=p\cdot x(p',w')\le w=p\cdot x(p,w). \]
Hence
\[ p\cdot\Bigl(x\Bigl(\frac{w}{w'}p',w\Bigr)-x(p,w)\Bigr)\le 0. \]
From the ULD (same wealth $w$, prices $\tfrac{w}{w'}p'$ and $p$; strict because
$x(\tfrac{w}{w'}p',w)=x(p',w')\ne x(p,w)$):
\[ \Bigl(\frac{w}{w'}p'-p\Bigr)\cdot\Bigl(x\Bigl(\frac{w}{w'}p',w\Bigr)-x(p,w)\Bigr)<0, \]
\[ \frac{w}{w'}p'\cdot\Bigl(x\Bigl(\frac{w}{w'}p',w\Bigr)-x(p,w)\Bigr)
   < \underbrace{p\cdot\Bigl(x\Bigl(\frac{w}{w'}p',w\Bigr)-x(p,w)\Bigr)}_{\le 0}. \]
We collect (dropping the positive factor $w/w'$):
\[ p'\cdot\Bigl(x\Bigl(\frac{w}{w'}p',w\Bigr)-x(p,w)\Bigr)<0, \qquad
   p'\cdot x\Bigl(\frac{w}{w'}p',w\Bigr)<p'\cdot x(p,w), \]
\[ p'\cdot x(p',w')<p'\cdot x(p,w)\quad\Longleftrightarrow\quad w'<p'\cdot x(p,w), \]
using Walras' law $p'\cdot x(p',w')=w'$. And this is WARP. \hfill$\square$

\topic{Where this leaves us}
We have seen
\[ \text{ULD}\ \Longrightarrow\ \text{WARP}\ \Longleftrightarrow\ \text{CLD}. \]
As already remarked, the ULD is stronger than the CLD.

\runin{Note.} It is enough that the ULD holds in the aggregate, even if it is
violated for some subset of agents, as long as the behavior of the remaining
agents offsets the behavior of the former.

However, whether this applies may depend on the wealth distribution rule. In our
example (Lesson 11), $\alpha_1=\alpha_2=\tfrac12$: this wealth distribution rule
was not enough for WARP, so the aggregate ULD must be violated. Change the rule to
$\alpha_1=1$, $\alpha_2=0$ $\Rightarrow$ aggregate demand $=$ agent 1's demand,
which satisfies WARP.

\topic{So far}
\begin{flatnum}
  \item Aggregate demand $\sum_i x_i(p,w_i)$ can be written as
        $x\bigl(p,\sum_i w_i\bigr)$ only in special cases. One: there is a
        deterministic wealth distribution rule, $w_i=w_i(p,w)$.
  \item If $w_i=\alpha_i w$, then aggregate demand can be written $x(p,w)$, and it
        satisfies WARP if individual demand satisfies the ULD.
\end{flatnum}


\section{Question 3: representative agents}

When does aggregate demand behave as if there were one agent with a preference
relation $\succsim$ that gave rise to the aggregate demand?

Such an agent is a \emph{positive} representative agent, because it describes
behavior.

However, even if there is a positive representative agent, it does not follow that
we can make welfare comparisons based on his preferences. Ideally, we would want a
\emph{normative} representative agent that we can use for welfare evaluation.

\topic{A social welfare function}
We define a social welfare function
\[ W(v_1,v_2,\dots,v_I) \]
that aggregates the utilities of all agents in the economy and gives a welfare
evaluation. It captures how trade-offs between agents are made as a society.

For now, take $W(v_1,\dots,v_I)$ as given, assuming it is increasing and concave.

\topic{Maximizing welfare}
We would like to maximize welfare. Given $(p,w)$, we can redistribute wealth across
agents $\Rightarrow$ in other words, the wealth distribution rule should be tailored
to maximize welfare: the distribution is its own choice variable. More formally,
solve
\[ \max_{w_1,w_2,\dots,w_I}\ W\bigl(v_1(p,w_1),v_2(p,w_2),\dots,v_I(p,w_I)\bigr)
   \qquad\text{s.t.}\quad \sum_{i=1}^{I}w_i\le w. \]
We can solve each agent's utility maximization problem to get $v_i(p,w_i)$, but
also the individual demand $x_i(p,w_i)$ $\Rightarrow$ we can calculate aggregate
demand.

\topic{The value function as an indirect utility}
Construct the maximum value function of this problem, i.e.\ welfare at the solution.
It depends only on the parameters $(p,w)$. Denote it $v(p,w)$. It can be shown that
$v(p,w)$ has all the usual properties of an indirect utility function, and that the
demand stemming from this indirect utility function is exactly the aggregate demand
we obtained before, $x(p,w)=\sum_i x_i\bigl(p,w_i(p,w)\bigr)$ with $w_i(p,w)$ the
optimal rule.

$\Rightarrow$ It is as if we have a normative representative agent with indirect
utility function $v(p,w)$. We can use it to derive aggregate demand and to evaluate
welfare.

This only requires an optimal wealth distribution rule.


\begin{transcriptnotes}
  \item Page 1 is a cover page reading ``Lesson 13'' only; the subtitle is written to
        match Lesson 11. No date is given. Lesson 12 (comparative statics) has slides
        and handouts only, no handwritten export, so it has no LaTeX lesson.
  \item Notation follows Lessons 7--11: $P$, $W$, $X$, $V$ are typeset $p$, $w$, $x$,
        $v$; the welfare function is $W$.
  \item Page 2: \textbf{Corrected:} the heading reads ``The compensated law of demand
        (CLD)'', but the definition that follows (income held fixed, $x_i(p',w_i)$
        vs $x_i(p,w_i)$) and every later page call it the \emph{uncompensated} law
        of demand (ULD); the heading is typeset ULD. \textbf{Added:} ``for any $p$,
        $p'$, $w_i$'' and the strict inequality when the bundles differ --- it is
        needed for the strict WARP conclusion (see page 4). ``satiiesfies'',
        ``satisfies CLD if'' $\to$ \emph{satisfies}; ``rules out some preferences it
        receives out Giffen goods'' $\to$ \emph{it rules out Giffen goods}.
        ``graphical example'' refers to the $\alpha_1=\alpha_2=\tfrac12$
        counterexample in Lesson 11; the pointer is \textbf{added}.
  \item Page 3: ``homothetic preferences satisfies'' $\to$ \emph{satisfy};
        ``aggegation ULD'' $\to$ \emph{the aggregate ULD}; ``in part to individual
        how homogeneity'' $\to$ \emph{in part on individual homogeneity}. The summed
        inequality in the parenthesis is \textbf{added}. ``We must show
        $p'\cdot x(p,w)>w'$'' is \textbf{added}.
  \item Page 4: ``homogenores'' $\to$ \emph{homogeneous}; ``individual demand are''
        $\to$ \emph{demands are}. The page writes, next to $p\cdot x(p',w')\le w$,
        ``correction this should be strict, correct it!''. The inequality it points
        at is the WARP hypothesis and stays weak; what must be strict is the ULD for
        different bundles, and with it every later step. \textbf{Corrected:} the
        ULD line and the next three inequalities are made strict ($<$), and the
        under-brace $\le 0$ is kept. ``By R.P.'' is typeset as ``by the
        revealed-preference assumption and Walras' law'' ($w=p\cdot x(p,w)$ is
        Walras' law). \textbf{Corrected:} ``we collect
        $p'(x(\tfrac{w}{w'}p',w)-x(p',w))\le 0$'' has $x(p',w)$ for $x(p,w)$.
  \item Page 5: \textbf{Corrected:} $p'\cdot x(p',w')\le p'x(p,w)\Leftrightarrow
        w'\le p'x(p,w)$ is made strict --- WARP requires $p'\cdot x(p,w)>w'$.
        ``aggretate'' $\to$ \emph{aggregate}; ``even if it violate'' $\to$
        \emph{is violated}; ``remaining'' kept; ``Distri. rule'' $\to$
        \emph{distribution rule}. ``so ULD must be violated'' is read as the
        \emph{aggregate} ULD (each agent may satisfy it). ``$\alpha_1=\alpha_2$
        $=\tfrac12$'' is written as a small stacked note; read as
        $\alpha_1=\alpha_2=\tfrac12$.
  \item Page 6: ``can be written as $x(p,\sum w_i)$ once in special case. one:''
        $\to$ \emph{only in special cases. One:}. ``satisfes'' $\to$
        \emph{satisfies}; ``Aggregate demand be written'' $\to$ \emph{can be
        written}; the two clauses of item 2 are reordered. ``with a preferences
        relation'' $\to$ \emph{a preference relation}; ``aggreagato'' $\to$
        \emph{aggregate}.
  \item Page 7: ``welfare companions'' $\to$ \emph{welfare comparisons}; ``nomative''
        $\to$ \emph{normative}; ``gives it a welfare evaluation'' $\to$ \emph{gives
        a welfare evaluation}; $W(v_1,v_2,v_I)$ $\to$ $W(v_1,v_2,\dots,v_I)$.
  \item Page 8: ``assuming $\uparrow$ and concave'' $\to$ \emph{increasing and
        concave}; ``It's own choice variable'' read as \emph{the distribution is its
        own choice variable}. \textbf{Corrected:} the objective is written
        $W(v_1(p,w),v_2(p,w),\dots,v_I(p,w_I))$; each agent's utility depends on
        her own wealth, $v_i(p,w_i)$. ``$\max$'' is \textbf{added} in front of the
        choice variables $w_1,\dots,w_I$. ``ability max problem'' $\to$
        \emph{utility maximization problem}.
  \item Page 9: ``maxmum'' $\to$ \emph{maximum}; ``weffare'' $\to$ \emph{welfare};
        ``reprsent'' $\to$ \emph{representative}; ``predit/derive'' (one word
        overwritten on the other) typeset \emph{derive}; ``and to welfare'' $\to$
        \emph{and to evaluate welfare}. The page ends the sentence ``\ldots is exactly
        the aggregate demand we obtained before?'' with a question mark; the answer
        is yes (MWG Proposition 4.D.1), so it is typeset as a statement, with the
        formula for aggregate demand under the optimal rule \textbf{added}.
  \item Page 10 holds the single line ``This only requires an optimal wealth distri.
        rule.''
  \item Pages of the original, for tracing back: page 1 the cover; page 2
        Question 2 and the ULD; page 3 homothetic preferences, $w_i=\alpha_i w$, the
        ULD aggregates, start of the proof; page 4 the proof; page 5 the end of the
        proof, ULD $\Rightarrow$ WARP $\Leftrightarrow$ CLD, the note on offsetting
        agents; page 6 ``so far'' and Question 3; page 7 positive vs normative
        agents and the social welfare function; page 8 the welfare-maximization
        problem; page 9 the value function $v(p,w)$; page 10 the last line.
\end{transcriptnotes}

\end{document}
